\section{问题重述}

\subsection{问题背景}

中药材采收后含水量高，必须及时干燥到安全水分以下才能贮藏与流通。热风烘干是产地初加工中
最常用的方式：热空气向表面供热并带走析出水分，水分则由芯部向表面扩散。热量传播快、水分
迁移慢，芯部往往在表面明显干燥之后很久才达到安全水分。干燥不足有霉变风险，过度延长则增加
能耗并可能损失热敏性有效成分\upcite{Rocha2011}。真实药材失水时还会收缩\upcite{Mayor2004}，且密度、比热容、
导热系数与水分扩散系数随含水率变化，扩散系数还依赖温度，传热与传质因此必须联立求解。
事先预测内部温度与水分的分布并给出可靠烘干时长，直接决定工艺参数与设备排产。

\subsection{问题的提出}

本题以圆柱形药材为对象。药材初始半径 $2$~cm、长度 $25$~cm，初始温度 $28~^\circ$C、
初始干基含水率 $2.55$~kg/kg，被送入烘房干燥。烘房内热空气的温度与含湿量随时间变化，
由附件给出的实测时序描述；药材表面与空气之间的对流换热系数和对流传质系数为已知常数。
四个子问题沿同一物理过程逐步放开限制条件：前两问要求刻画给定时段内的场分布，
后两问要求给出达到安全水分所需的烘干时长，其中最后一问还需考虑药材的实际收缩。

问题一考察预热阶段。此时物性可按常数处理，水分扩散系数只随局部含水率变化。
需要建立药材内部的温度与水分输运关系，求出前 $1800$~s 内若干时刻沿半径若干位置的
温度与含水率，并给出全过程的完整场数据。

问题二把物性的真实依赖关系全部放开。密度、比热容与导热系数随含水率变化，水分扩散系数
同时随含水率和温度变化，温度场与水分场因此互相牵动。需要在这一耦合条件下重新刻画整个
烘干过程，并给出前 $3$~h 内的温度与含水率分布。

问题三在问题二的模型上追加干燥完成的判定。当药材内部含水率降到 $0.15$~kg/kg 以下即视为
烘干完成，需要给出达到该状态所需的时长，并给出全过程中若干时刻的含水率分布。

问题四进一步考虑药材的径向收缩。半径随时间的实测数据由附件给出，物性关系也相应更换。
求解域不再固定，需要在随时间缩小的区域上重建模型，重新给出烘干时长与含水率分布，
并说明收缩这一因素究竟把烘干时长改变了多少。
